% Clinician review packet -- GENERATED by scripts/build_clinician_review.py.
% Editable/fillable PDF: compile with pdflatex (the checkboxes and text fields are
% AcroForm fields via hyperref; open the PDF in any reader that supports forms).
\documentclass[10pt,a4paper]{article}
\usepackage[T1]{fontenc}
\usepackage[utf8]{inputenc}
\usepackage[a4paper,margin=1.5cm,top=1.6cm,bottom=1.6cm]{geometry}
\usepackage{longtable}
\usepackage{array}
\usepackage{xcolor}
\usepackage{colortbl}
\usepackage{amssymb}
\usepackage{enumitem}
\usepackage[colorlinks=false,pdfborder={0 0 0}]{hyperref}
\usepackage{titlesec}

\definecolor{refer}{HTML}{C2255C}
\definecolor{treat}{HTML}{B45309}
\definecolor{home}{HTML}{2F8F4E}
\definecolor{accent}{HTML}{0E7490}
\definecolor{warnbg}{HTML}{FBEEE1}
\definecolor{rule}{HTML}{DDE5EA}
\newcommand{\refer}[1]{\textcolor{refer}{\textbf{#1}}}
\newcommand{\treat}[1]{\textcolor{treat}{\textbf{#1}}}
\newcommand{\home}[1]{\textcolor{home}{\textbf{#1}}}
\titleformat{\section}{\large\bfseries\color{accent}}{}{0em}{}
\titleformat{\subsection}{\normalsize\bfseries}{}{0em}{}
\setlength{\parindent}{0pt}
\setlength{\extrarowheight}{2pt}
\renewcommand{\arraystretch}{1.15}
\hypersetup{pdftitle={IMCI extended + young-infant branches -- clinician review}}

\begin{document}
\begin{Form}

{\Large\bfseries\color{accent} IMCI extended \& young-infant branches --- clinician review}\\[2pt]
{\footnotesize\color{accent} IMCI \textperiodcentered{} clinical sign-off}\\[8pt]
\colorbox{warnbg}{\parbox{\dimexpr\linewidth-2\fboxsep}{\textbf{\color{treat}UNREVIEWED.} These classifications and dose bands train a model only with \texttt{-{}-include-extended} (off by default). Nothing ships until signed off against the current national IMCI adaptation. Every row here is generated from the code, so it matches exactly what the model was trained to say.}}\\[8pt]
{\footnotesize Sources: \texttt{data/\allowbreak imci\_2022/\allowbreak classifications.json} (2022 SA adaptation, cross-checked vs the WHO 2014 generic); doses in \texttt{data/\allowbreak imci\_2022/\allowbreak dosing\_tables.json}. Severity uses the IMCI colour tiers. To review: confirm each \textbf{trigger} and \textbf{severity}, tick \textbf{OK} or write a correction, then sign at the end.}\\[6pt]
\section*{Classifications --- severity, trigger \& action (as implemented)}
\subsection*{Fever / malaria}
\begin{longtable}{@{}>{\ttfamily\footnotesize\raggedright\arraybackslash}p{2.7cm} >{\footnotesize\raggedright\arraybackslash}p{1.7cm} >{\footnotesize\raggedright\arraybackslash}p{4.2cm} >{\footnotesize\raggedright\arraybackslash}p{4.2cm} >{\footnotesize\raggedright\arraybackslash}p{1.9cm} >{\centering\arraybackslash}p{0.7cm}@{}}
\rowcolor{rule!40}\normalfont\footnotesize\textbf{Classification} & \textbf{Severity} & \textbf{Trigger (as coded)} & \textbf{Action} & \textbf{Drugs} & \textbf{OK}\\
\endhead
very\_\allowbreak severe\_\allowbreak febrile\_\allowbreak disease & \refer{PINK -- refer} & Fever with a general danger sign or a stiff neck classifies as very severe febrile disease, whatever the malaria test shows. & Give the first dose of an appropriate antibiotic. Treat the child to prevent low blood sugar. Give paracetamol for high fever. Refer URGENTLY to hospital. & ceftriaxone\_\allowbreak im, paracetamol & \CheckBox[name=ok_very_severe_febrile_disease_1,width=1.6ex,height=1.6ex]{}\\
malaria & \treat{YELLOW -- treat} & The malaria test is positive and there are no danger signs or stiff neck. & Give a first-line oral antimalarial as the national guideline directs. Give paracetamol for high fever. Advise the mother when to return immediately. Follow up in 3 days if the fever persists. & artemether\_\allowbreak lumefantrine, paracetamol & \CheckBox[name=ok_malaria_2,width=1.6ex,height=1.6ex]{}\\
fever\_\allowbreak no\_\allowbreak malaria & \home{GREEN -- home} & The malaria test is negative, so this fever is not malaria. & Look for and treat another cause of fever. Give paracetamol for high fever. Advise the mother when to return immediately. Follow up in 3 days if the fever persists. & paracetamol & \CheckBox[name=ok_fever_no_malaria_3,width=1.6ex,height=1.6ex]{}\\
fever\_\allowbreak malaria\_\allowbreak test\_\allowbreak required & \treat{YELLOW -- treat} & This child has fever in a malaria risk area, or has travelled to one, and no malaria test has been done. The chart booklet requires a malaria test before the fever can be classified. & Do a malaria test (RDT or microscopy) now and classify on the result. Give paracetamol for high fever meanwhile, and refer urgently if any danger sign appears. & paracetamol & \CheckBox[name=ok_fever_malaria_test_required_4,width=1.6ex,height=1.6ex]{}\\
\end{longtable}
\vspace{3pt}
{\footnotesize\textbf{Corrections for Fever / malaria:}}\par\vspace{2pt}
\TextField[name=tf_corr_Fever___malaria_5,width=\linewidth,height=1.2cm,multiline=true,backgroundcolor={1 1 0.98}]{}
\vspace{9pt}
\subsection*{Measles}
\begin{longtable}{@{}>{\ttfamily\footnotesize\raggedright\arraybackslash}p{2.7cm} >{\footnotesize\raggedright\arraybackslash}p{1.7cm} >{\footnotesize\raggedright\arraybackslash}p{4.2cm} >{\footnotesize\raggedright\arraybackslash}p{4.2cm} >{\footnotesize\raggedright\arraybackslash}p{1.9cm} >{\centering\arraybackslash}p{0.7cm}@{}}
\rowcolor{rule!40}\normalfont\footnotesize\textbf{Classification} & \textbf{Severity} & \textbf{Trigger (as coded)} & \textbf{Action} & \textbf{Drugs} & \textbf{OK}\\
\endhead
severe\_\allowbreak complicated\_\allowbreak measles & \refer{PINK -- refer} & A generalised rash with cough, runny nose, or red eyes meets the measles case definition. There is clouding of the cornea, which makes this severe complicated measles. & Give vitamin A. Give the first dose of an appropriate antibiotic. If there is clouding of the cornea or pus draining from the eye, apply tetracycline eye ointment. Refer URGENTLY to hospital. & vitamin\_\allowbreak a & \CheckBox[name=ok_severe_complicated_measles_6,width=1.6ex,height=1.6ex]{}\\
measles\_\allowbreak with\_\allowbreak eye\_\allowbreak or\_\allowbreak mouth\_\allowbreak complications & \treat{YELLOW -- treat} & A generalised rash with cough, runny nose, or red eyes meets the measles case definition. There is pus draining from the eye. & Give vitamin A. If there is pus draining from the eye, apply tetracycline eye ointment. If there are mouth ulcers, treat with gentian violet. Follow up in 3 days. & vitamin\_\allowbreak a & \CheckBox[name=ok_measles_with_eye_or_mouth_complications_7,width=1.6ex,height=1.6ex]{}\\
measles & \home{GREEN -- home} & A generalised rash with cough, runny nose, or red eyes meets the measles case definition. There are no eye or mouth complications and no danger signs. & Give vitamin A. Advise the mother when to return immediately. & vitamin\_\allowbreak a & \CheckBox[name=ok_measles_8,width=1.6ex,height=1.6ex]{}\\
\end{longtable}
\vspace{3pt}
{\footnotesize\textbf{Corrections for Measles:}}\par\vspace{2pt}
\TextField[name=tf_corr_Measles_9,width=\linewidth,height=1.2cm,multiline=true,backgroundcolor={1 1 0.98}]{}
\vspace{9pt}
\subsection*{Anaemia}
\begin{longtable}{@{}>{\ttfamily\footnotesize\raggedright\arraybackslash}p{2.7cm} >{\footnotesize\raggedright\arraybackslash}p{1.7cm} >{\footnotesize\raggedright\arraybackslash}p{4.2cm} >{\footnotesize\raggedright\arraybackslash}p{4.2cm} >{\footnotesize\raggedright\arraybackslash}p{1.9cm} >{\centering\arraybackslash}p{0.7cm}@{}}
\rowcolor{rule!40}\normalfont\footnotesize\textbf{Classification} & \textbf{Severity} & \textbf{Trigger (as coded)} & \textbf{Action} & \textbf{Drugs} & \textbf{OK}\\
\endhead
severe\_\allowbreak anaemia & \refer{PINK -- refer} & There is severe palmar pallor. & Refer URGENTLY to hospital. & -- & \CheckBox[name=ok_severe_anaemia_10,width=1.6ex,height=1.6ex]{}\\
anaemia & \treat{YELLOW -- treat} & There is some palmar pallor. & Give iron. Give mebendazole if the child is 1 year or older and has not had a dose in the last 6 months. Advise the mother when to return immediately. Follow up in 14 days. & -- & \CheckBox[name=ok_anaemia_11,width=1.6ex,height=1.6ex]{}\\
no\_\allowbreak anaemia & \home{GREEN -- home} & There is no palmar pallor. & No treatment for anaemia is needed. & -- & \CheckBox[name=ok_no_anaemia_12,width=1.6ex,height=1.6ex]{}\\
\end{longtable}
\vspace{3pt}
{\footnotesize\textbf{Corrections for Anaemia:}}\par\vspace{2pt}
\TextField[name=tf_corr_Anaemia_13,width=\linewidth,height=1.2cm,multiline=true,backgroundcolor={1 1 0.98}]{}
\vspace{9pt}
\subsection*{Acute malnutrition}
\begin{longtable}{@{}>{\ttfamily\footnotesize\raggedright\arraybackslash}p{2.7cm} >{\footnotesize\raggedright\arraybackslash}p{1.7cm} >{\footnotesize\raggedright\arraybackslash}p{4.2cm} >{\footnotesize\raggedright\arraybackslash}p{4.2cm} >{\footnotesize\raggedright\arraybackslash}p{1.9cm} >{\centering\arraybackslash}p{0.7cm}@{}}
\rowcolor{rule!40}\normalfont\footnotesize\textbf{Classification} & \textbf{Severity} & \textbf{Trigger (as coded)} & \textbf{Action} & \textbf{Drugs} & \textbf{OK}\\
\endhead
complicated\_\allowbreak severe\_\allowbreak acute\_\allowbreak malnutrition & \refer{PINK -- refer} & There is oedema of both feet, which is always severe acute malnutrition. With oedema of both feet, this is complicated severe acute malnutrition. & Give the first dose of an appropriate antibiotic. Treat the child to prevent low blood sugar. Keep the child warm. Refer URGENTLY to hospital. & vitamin\_\allowbreak a & \CheckBox[name=ok_complicated_severe_acute_malnutrition_14,width=1.6ex,height=1.6ex]{}\\
uncomplicated\_\allowbreak severe\_\allowbreak acute\_\allowbreak malnutrition & \treat{YELLOW -- treat} & MUAC is 110mm, below 115mm. The appetite test passed and there is no medical complication. & Give ready-to-use therapeutic food (RUTF) for the child to take at home. Assess and counsel on feeding. Advise the mother when to return immediately. Follow up in 7 days. & amoxicillin, vitamin\_\allowbreak a & \CheckBox[name=ok_uncomplicated_severe_acute_malnutrition_15,width=1.6ex,height=1.6ex]{}\\
moderate\_\allowbreak acute\_\allowbreak malnutrition & \treat{YELLOW -- treat} & MUAC is 120mm, between 115mm and 125mm. & Assess and counsel on feeding. Advise the mother when to return immediately. Follow up in 30 days. & vitamin\_\allowbreak a & \CheckBox[name=ok_moderate_acute_malnutrition_16,width=1.6ex,height=1.6ex]{}\\
no\_\allowbreak acute\_\allowbreak malnutrition & \home{GREEN -- home} & There is no oedema of both feet and no wasting by MUAC or weight-for-height. & If the child is less than 2 years old, assess and counsel on feeding. No treatment for acute malnutrition is needed. & -- & \CheckBox[name=ok_no_acute_malnutrition_17,width=1.6ex,height=1.6ex]{}\\
\end{longtable}
\vspace{3pt}
{\footnotesize\textbf{Corrections for Acute malnutrition:}}\par\vspace{2pt}
\TextField[name=tf_corr_Acute_malnutrition_18,width=\linewidth,height=1.2cm,multiline=true,backgroundcolor={1 1 0.98}]{}
\vspace{9pt}
\subsection*{Wheeze (cough sub-branch)}
\begin{longtable}{@{}>{\ttfamily\footnotesize\raggedright\arraybackslash}p{2.7cm} >{\footnotesize\raggedright\arraybackslash}p{1.7cm} >{\footnotesize\raggedright\arraybackslash}p{4.2cm} >{\footnotesize\raggedright\arraybackslash}p{4.2cm} >{\footnotesize\raggedright\arraybackslash}p{1.9cm} >{\centering\arraybackslash}p{0.7cm}@{}}
\rowcolor{rule!40}\normalfont\footnotesize\textbf{Classification} & \textbf{Severity} & \textbf{Trigger (as coded)} & \textbf{Action} & \textbf{Drugs} & \textbf{OK}\\
\endhead
wheeze\_\allowbreak with\_\allowbreak danger\_\allowbreak sign & \refer{PINK -- refer} & The child has wheeze. There is a general danger sign, so give salbutamol and refer urgently. & Give salbutamol by spacer. Give the first dose of an appropriate antibiotic and refer URGENTLY to hospital. & -- & \CheckBox[name=ok_wheeze_with_danger_sign_19,width=1.6ex,height=1.6ex]{}\\
wheeze & \treat{YELLOW -- treat} & The child has wheeze. & Give salbutamol by spacer for 5 days. Follow up in 5 days if the child is still wheezing. If a cough lasts more than 14 days or the wheeze is recurrent, assess for TB or asthma. & -- & \CheckBox[name=ok_wheeze_20,width=1.6ex,height=1.6ex]{}\\
\end{longtable}
\vspace{3pt}
{\footnotesize\textbf{Corrections for Wheeze (cough sub-branch):}}\par\vspace{2pt}
\TextField[name=tf_corr_Wheeze__cough_sub_branch__21,width=\linewidth,height=1.2cm,multiline=true,backgroundcolor={1 1 0.98}]{}
\vspace{9pt}
\subsection*{Persistent diarrhoea (>=14 days)}
\begin{longtable}{@{}>{\ttfamily\footnotesize\raggedright\arraybackslash}p{2.7cm} >{\footnotesize\raggedright\arraybackslash}p{1.7cm} >{\footnotesize\raggedright\arraybackslash}p{4.2cm} >{\footnotesize\raggedright\arraybackslash}p{4.2cm} >{\footnotesize\raggedright\arraybackslash}p{1.9cm} >{\centering\arraybackslash}p{0.7cm}@{}}
\rowcolor{rule!40}\normalfont\footnotesize\textbf{Classification} & \textbf{Severity} & \textbf{Trigger (as coded)} & \textbf{Action} & \textbf{Drugs} & \textbf{OK}\\
\endhead
severe\_\allowbreak persistent\_\allowbreak diarrhoea & \refer{PINK -- refer} & Diarrhoea has lasted 14 days or more and dehydration is present. & Treat dehydration before referral unless the child has another severe classification. Give an extra dose of vitamin A. Refer URGENTLY to hospital. & vitamin\_\allowbreak a & \CheckBox[name=ok_severe_persistent_diarrhoea_22,width=1.6ex,height=1.6ex]{}\\
persistent\_\allowbreak diarrhoea & \treat{YELLOW -- treat} & Diarrhoea has lasted 14 days or more with no visible dehydration. & Counsel the caregiver on feeding for persistent diarrhoea. Give an extra dose of vitamin A and give zinc for 14 days. Follow up in 5 days. & zinc, vitamin\_\allowbreak a & \CheckBox[name=ok_persistent_diarrhoea_23,width=1.6ex,height=1.6ex]{}\\
\end{longtable}
\vspace{3pt}
{\footnotesize\textbf{Corrections for Persistent diarrhoea (>=14 days):}}\par\vspace{2pt}
\TextField[name=tf_corr_Persistent_diarrhoea____14_days__24,width=\linewidth,height=1.2cm,multiline=true,backgroundcolor={1 1 0.98}]{}
\vspace{9pt}
\subsection*{Dysentery (blood in stool)}
\begin{longtable}{@{}>{\ttfamily\footnotesize\raggedright\arraybackslash}p{2.7cm} >{\footnotesize\raggedright\arraybackslash}p{1.7cm} >{\footnotesize\raggedright\arraybackslash}p{4.2cm} >{\footnotesize\raggedright\arraybackslash}p{4.2cm} >{\footnotesize\raggedright\arraybackslash}p{1.9cm} >{\centering\arraybackslash}p{0.7cm}@{}}
\rowcolor{rule!40}\normalfont\footnotesize\textbf{Classification} & \textbf{Severity} & \textbf{Trigger (as coded)} & \textbf{Action} & \textbf{Drugs} & \textbf{OK}\\
\endhead
severe\_\allowbreak dysentery & \refer{PINK -- refer} & There is blood in the stool and the child is less than 12 months old. & Treat the child to prevent low blood sugar, keep the child warm, and refer URGENTLY to hospital. & -- & \CheckBox[name=ok_severe_dysentery_25,width=1.6ex,height=1.6ex]{}\\
dysentery & \treat{YELLOW -- treat} & There is blood in the stool, the child is 12 months or older, and there is no dehydration. & Treat for 3 days with ciprofloxacin. Advise the mother when to return immediately. Follow up in 2 days. & -- & \CheckBox[name=ok_dysentery_26,width=1.6ex,height=1.6ex]{}\\
\end{longtable}
\vspace{3pt}
{\footnotesize\textbf{Corrections for Dysentery (blood in stool):}}\par\vspace{2pt}
\TextField[name=tf_corr_Dysentery__blood_in_stool__27,width=\linewidth,height=1.2cm,multiline=true,backgroundcolor={1 1 0.98}]{}
\vspace{9pt}
\subsection*{Sore throat (from 3 years)}
\begin{longtable}{@{}>{\ttfamily\footnotesize\raggedright\arraybackslash}p{2.7cm} >{\footnotesize\raggedright\arraybackslash}p{1.7cm} >{\footnotesize\raggedright\arraybackslash}p{4.2cm} >{\footnotesize\raggedright\arraybackslash}p{4.2cm} >{\footnotesize\raggedright\arraybackslash}p{1.9cm} >{\centering\arraybackslash}p{0.7cm}@{}}
\rowcolor{rule!40}\normalfont\footnotesize\textbf{Classification} & \textbf{Severity} & \textbf{Trigger (as coded)} & \textbf{Action} & \textbf{Drugs} & \textbf{OK}\\
\endhead
streptococcal\_\allowbreak sore\_\allowbreak throat & \treat{YELLOW -- treat} & There is enlarged tonsils, white or yellow exudate on the tonsils, with no runny nose and no cough, which points to a streptococcal sore throat. & Give penicillin. Treat pain and fever. Soothe the throat with a safe remedy. Follow up in 5 days if symptoms are worse or not resolving. & -- & \CheckBox[name=ok_streptococcal_sore_throat_28,width=1.6ex,height=1.6ex]{}\\
sore\_\allowbreak throat\_\allowbreak non\_\allowbreak streptococcal & \home{GREEN -- home} & There are not enough signs to classify this as a streptococcal sore throat. & Soothe the throat with a safe remedy. & -- & \CheckBox[name=ok_sore_throat_non_streptococcal_29,width=1.6ex,height=1.6ex]{}\\
\end{longtable}
\vspace{3pt}
{\footnotesize\textbf{Corrections for Sore throat (from 3 years):}}\par\vspace{2pt}
\TextField[name=tf_corr_Sore_throat__from_3_years__30,width=\linewidth,height=1.2cm,multiline=true,backgroundcolor={1 1 0.98}]{}
\vspace{9pt}
\subsection*{Growth problem (RTHB curve)}
\begin{longtable}{@{}>{\ttfamily\footnotesize\raggedright\arraybackslash}p{2.7cm} >{\footnotesize\raggedright\arraybackslash}p{1.7cm} >{\footnotesize\raggedright\arraybackslash}p{4.2cm} >{\footnotesize\raggedright\arraybackslash}p{4.2cm} >{\footnotesize\raggedright\arraybackslash}p{1.9cm} >{\centering\arraybackslash}p{0.7cm}@{}}
\rowcolor{rule!40}\normalfont\footnotesize\textbf{Classification} & \textbf{Severity} & \textbf{Trigger (as coded)} & \textbf{Action} & \textbf{Drugs} & \textbf{OK}\\
\endhead
growth\_\allowbreak problem & \treat{YELLOW -- treat} & On the weight curve, the child is losing weight. & Assess feeding and counsel the caregiver on the feeding recommendations. Deworm and give vitamin A if due. Advise the mother when to return immediately. Follow up in 7 days if there is a feeding problem, otherwise in 14 days. & -- & \CheckBox[name=ok_growth_problem_31,width=1.6ex,height=1.6ex]{}\\
\end{longtable}
\vspace{3pt}
{\footnotesize\textbf{Corrections for Growth problem (RTHB curve):}}\par\vspace{2pt}
\TextField[name=tf_corr_Growth_problem__RTHB_curve__32,width=\linewidth,height=1.2cm,multiline=true,backgroundcolor={1 1 0.98}]{}
\vspace{9pt}
\subsection*{HIV}
\begin{longtable}{@{}>{\ttfamily\footnotesize\raggedright\arraybackslash}p{2.7cm} >{\footnotesize\raggedright\arraybackslash}p{1.7cm} >{\footnotesize\raggedright\arraybackslash}p{4.2cm} >{\footnotesize\raggedright\arraybackslash}p{4.2cm} >{\footnotesize\raggedright\arraybackslash}p{1.9cm} >{\centering\arraybackslash}p{0.7cm}@{}}
\rowcolor{rule!40}\normalfont\footnotesize\textbf{Classification} & \textbf{Severity} & \textbf{Trigger (as coded)} & \textbf{Action} & \textbf{Drugs} & \textbf{OK}\\
\endhead
confirmed\_\allowbreak hiv\_\allowbreak infection & \treat{YELLOW -- treat} & The child has a positive HIV test or is already on ART, so HIV infection is confirmed. & Follow the steps to initiate ART. Give cotrimoxazole prophylaxis from 6 weeks of age. Ask about the caregiver's health and manage appropriately. Provide long-term follow-up. ART initiation is not urgent -- stabilise any severe illness first. & -- & \CheckBox[name=ok_confirmed_hiv_infection_33,width=1.6ex,height=1.6ex]{}\\
hiv\_\allowbreak exposed & \treat{YELLOW -- treat} & The child is HIV-exposed: on ARV prophylaxis, or with a negative test while still breastfeeding or within 6 weeks of breastfeeding, so infection is not yet ruled out. & Complete the appropriate infant ARV prophylaxis. Repeat HIV PCR testing per the schedule and reclassify on the result. Ask about the caregiver's health and provide follow-up care. & -- & \CheckBox[name=ok_hiv_exposed_34,width=1.6ex,height=1.6ex]{}\\
suspected\_\allowbreak symptomatic\_\allowbreak hiv & \treat{YELLOW -- treat} & There are 3 features of HIV infection (three or more), so symptomatic HIV is suspected. & Counsel and offer HIV testing for the child and reclassify on the result. Counsel the caregiver about her own health and offer testing. Provide long-term follow-up. & -- & \CheckBox[name=ok_suspected_symptomatic_hiv_35,width=1.6ex,height=1.6ex]{}\\
possible\_\allowbreak hiv\_\allowbreak infection & \treat{YELLOW -- treat} & HIV infection is possible because the mother is HIV-positive. & Provide routine care including HIV testing for the child. Counsel the caregiver about her health and offer testing and treatment as needed. Reclassify on the test result. & -- & \CheckBox[name=ok_possible_hiv_infection_36,width=1.6ex,height=1.6ex]{}\\
hiv\_\allowbreak infection\_\allowbreak unlikely & \home{GREEN -- home} & The HIV test is negative, all breastfeeding stopped 6 weeks or more before the test, and there are no features of HIV infection. & Provide routine care. Repeat HIV testing only if new features appear or exposure is ongoing. & -- & \CheckBox[name=ok_hiv_infection_unlikely_37,width=1.6ex,height=1.6ex]{}\\
\end{longtable}
\vspace{3pt}
{\footnotesize\textbf{Corrections for HIV:}}\par\vspace{2pt}
\TextField[name=tf_corr_HIV_38,width=\linewidth,height=1.2cm,multiline=true,backgroundcolor={1 1 0.98}]{}
\vspace{9pt}
\subsection*{Young infant: bacterial infection}
\begin{longtable}{@{}>{\ttfamily\footnotesize\raggedright\arraybackslash}p{2.7cm} >{\footnotesize\raggedright\arraybackslash}p{1.7cm} >{\footnotesize\raggedright\arraybackslash}p{4.2cm} >{\footnotesize\raggedright\arraybackslash}p{4.2cm} >{\footnotesize\raggedright\arraybackslash}p{1.9cm} >{\centering\arraybackslash}p{0.7cm}@{}}
\rowcolor{rule!40}\normalfont\footnotesize\textbf{Classification} & \textbf{Severity} & \textbf{Trigger (as coded)} & \textbf{Action} & \textbf{Drugs} & \textbf{OK}\\
\endhead
yi\_\allowbreak very\_\allowbreak severe\_\allowbreak disease & \refer{PINK -- refer} & The young infant has a bulging fontanelle, which is very severe disease. & Give the first dose of an appropriate intramuscular antibiotic. Treat to prevent low blood sugar. Keep the infant warm on the way to hospital. Refer URGENTLY. & -- & \CheckBox[name=ok_yi_very_severe_disease_39,width=1.6ex,height=1.6ex]{}\\
yi\_\allowbreak local\_\allowbreak bacterial\_\allowbreak infection & \treat{YELLOW -- treat} & The young infant has a red umbilicus, a local bacterial infection. & Give an appropriate oral antibiotic. If there is eye discharge, give an eye ointment. Teach the caregiver to treat the local infection at home and to give home care. Follow up in 2 days. & -- & \CheckBox[name=ok_yi_local_bacterial_infection_40,width=1.6ex,height=1.6ex]{}\\
yi\_\allowbreak no\_\allowbreak bacterial\_\allowbreak infection & \home{GREEN -- home} & No signs of very severe disease or a local bacterial infection. & Counsel the caregiver on home care for the young infant. & -- & \CheckBox[name=ok_yi_no_bacterial_infection_41,width=1.6ex,height=1.6ex]{}\\
\end{longtable}
\vspace{3pt}
{\footnotesize\textbf{Corrections for Young infant: bacterial infection:}}\par\vspace{2pt}
\TextField[name=tf_corr_Young_infant__bacterial_infection_42,width=\linewidth,height=1.2cm,multiline=true,backgroundcolor={1 1 0.98}]{}
\vspace{9pt}
\subsection*{Young infant: jaundice}
\begin{longtable}{@{}>{\ttfamily\footnotesize\raggedright\arraybackslash}p{2.7cm} >{\footnotesize\raggedright\arraybackslash}p{1.7cm} >{\footnotesize\raggedright\arraybackslash}p{4.2cm} >{\footnotesize\raggedright\arraybackslash}p{4.2cm} >{\footnotesize\raggedright\arraybackslash}p{1.9cm} >{\centering\arraybackslash}p{0.7cm}@{}}
\rowcolor{rule!40}\normalfont\footnotesize\textbf{Classification} & \textbf{Severity} & \textbf{Trigger (as coded)} & \textbf{Action} & \textbf{Drugs} & \textbf{OK}\\
\endhead
yi\_\allowbreak severe\_\allowbreak jaundice & \refer{PINK -- refer} & There is jaundice under 24 hours of age, which is severe jaundice. & Treat to prevent low blood sugar. Keep the infant warm. Refer URGENTLY to hospital. & -- & \CheckBox[name=ok_yi_severe_jaundice_43,width=1.6ex,height=1.6ex]{}\\
yi\_\allowbreak jaundice & \treat{YELLOW -- treat} & There is jaundice appearing after 24 hours of age, with palms and soles not yellow. & Advise the caregiver to return immediately if the palms and soles become yellow. Follow up in 1 day. If the infant is older than 14 days, refer for assessment. & -- & \CheckBox[name=ok_yi_jaundice_44,width=1.6ex,height=1.6ex]{}\\
\end{longtable}
\vspace{3pt}
{\footnotesize\textbf{Corrections for Young infant: jaundice:}}\par\vspace{2pt}
\TextField[name=tf_corr_Young_infant__jaundice_45,width=\linewidth,height=1.2cm,multiline=true,backgroundcolor={1 1 0.98}]{}
\vspace{9pt}
\subsection*{Young infant: diarrhoea}
\begin{longtable}{@{}>{\ttfamily\footnotesize\raggedright\arraybackslash}p{2.7cm} >{\footnotesize\raggedright\arraybackslash}p{1.7cm} >{\footnotesize\raggedright\arraybackslash}p{4.2cm} >{\footnotesize\raggedright\arraybackslash}p{4.2cm} >{\footnotesize\raggedright\arraybackslash}p{1.9cm} >{\centering\arraybackslash}p{0.7cm}@{}}
\rowcolor{rule!40}\normalfont\footnotesize\textbf{Classification} & \textbf{Severity} & \textbf{Trigger (as coded)} & \textbf{Action} & \textbf{Drugs} & \textbf{OK}\\
\endhead
yi\_\allowbreak severe\_\allowbreak dehydration & \refer{PINK -- refer} & There is diarrhoea with the infant is less than 1 month old, which is severe dehydration in a young infant. & Start intravenous fluids (Plan C). Give the first dose of an intramuscular antibiotic. Keep the infant warm on the way to hospital. Refer URGENTLY. & -- & \CheckBox[name=ok_yi_severe_dehydration_46,width=1.6ex,height=1.6ex]{}\\
yi\_\allowbreak dysentery & \refer{PINK -- refer} & There is blood in the stool of a young infant. & Keep the infant warm on the way to hospital. Refer URGENTLY. & -- & \CheckBox[name=ok_yi_dysentery_47,width=1.6ex,height=1.6ex]{}\\
yi\_\allowbreak severe\_\allowbreak persistent\_\allowbreak diarrhoea & \refer{PINK -- refer} & Diarrhoea has lasted 14 days or more in a young infant. & Treat any dehydration before referral. Keep the infant warm. Refer URGENTLY. & -- & \CheckBox[name=ok_yi_severe_persistent_diarrhoea_48,width=1.6ex,height=1.6ex]{}\\
yi\_\allowbreak some\_\allowbreak dehydration & \treat{YELLOW -- treat} & There is diarrhoea with two of restless/irritable, sunken eyes, and a slow skin pinch, which is some dehydration. & Give fluid for some dehydration (Plan B). Advise the mother to continue breastfeeding. Give zinc for 14 days. Follow up in 2 days. & -- & \CheckBox[name=ok_yi_some_dehydration_49,width=1.6ex,height=1.6ex]{}\\
yi\_\allowbreak no\_\allowbreak dehydration & \home{GREEN -- home} & There is diarrhoea with not enough signs for some or severe dehydration. & Give fluid and continue breastfeeding at home (Plan A). Give zinc for 14 days. Counsel the caregiver on home care. Follow up in 2 days. & -- & \CheckBox[name=ok_yi_no_dehydration_50,width=1.6ex,height=1.6ex]{}\\
\end{longtable}
\vspace{3pt}
{\footnotesize\textbf{Corrections for Young infant: diarrhoea:}}\par\vspace{2pt}
\TextField[name=tf_corr_Young_infant__diarrhoea_51,width=\linewidth,height=1.2cm,multiline=true,backgroundcolor={1 1 0.98}]{}
\vspace{9pt}
\subsection*{Young infant: congenital problems}
\begin{longtable}{@{}>{\ttfamily\footnotesize\raggedright\arraybackslash}p{2.7cm} >{\footnotesize\raggedright\arraybackslash}p{1.7cm} >{\footnotesize\raggedright\arraybackslash}p{4.2cm} >{\footnotesize\raggedright\arraybackslash}p{4.2cm} >{\footnotesize\raggedright\arraybackslash}p{1.9cm} >{\centering\arraybackslash}p{0.7cm}@{}}
\rowcolor{rule!40}\normalfont\footnotesize\textbf{Classification} & \textbf{Severity} & \textbf{Trigger (as coded)} & \textbf{Action} & \textbf{Drugs} & \textbf{OK}\\
\endhead
yi\_\allowbreak congenital\_\allowbreak priority & \refer{PINK -- refer} & There is a cleft lip or palate, a congenital priority sign. & Give any needed pre-referral treatment. Treat to prevent low blood sugar. Keep the infant warm. Refer URGENTLY to hospital. & -- & \CheckBox[name=ok_yi_congenital_priority_52,width=1.6ex,height=1.6ex]{}\\
yi\_\allowbreak congenital\_\allowbreak abnormal\_\allowbreak signs & \treat{YELLOW -- treat} & There is a club foot, an abnormal sign. & Keep the infant warm, skin to skin. Assess breastfeeding and support the mother. Refer for assessment. & -- & \CheckBox[name=ok_yi_congenital_abnormal_signs_53,width=1.6ex,height=1.6ex]{}\\
yi\_\allowbreak possible\_\allowbreak congenital\_\allowbreak syphilis & \treat{YELLOW -- treat} & The mother's RPR is positive and she was untreated or only partially treated, so congenital syphilis is possible. & Check for signs of congenital syphilis and refer to hospital if present. If there are no signs, give intramuscular penicillin. Ensure the mother receives full treatment. & -- & \CheckBox[name=ok_yi_possible_congenital_syphilis_54,width=1.6ex,height=1.6ex]{}\\
\end{longtable}
\vspace{3pt}
{\footnotesize\textbf{Corrections for Young infant: congenital problems:}}\par\vspace{2pt}
\TextField[name=tf_corr_Young_infant__congenital_problems_55,width=\linewidth,height=1.2cm,multiline=true,backgroundcolor={1 1 0.98}]{}
\vspace{9pt}
\section*{Dosing tables --- confirm every band}
{\footnotesize Doses are not graded by the model scorer; they are emitted verbatim from these tables.}\\[4pt]
\subsection*{amoxicillin}
{\footnotesize pneumonia, acute ear infection \textperiodcentered{} oral \textperiodcentered{} two times daily for 5 days \textperiodcentered{} source 2014 p16 \textperiodcentered{} keyed by weight}\\[-2pt]
\begin{longtable}{@{}l r r c@{}}
\rowcolor{rule!40}\footnotesize\textbf{weight (kg)} & \footnotesize\textbf{tablet\_250mg} & \footnotesize\textbf{syrup\_250mg\_per\_5ml\_ml} & \footnotesize\textbf{OK}\\
\endhead
\footnotesize 4--10 & \footnotesize 1 & \footnotesize 5 & \CheckBox[name=ok_amoxicillin_56,width=1.6ex,height=1.6ex]{}\\
\footnotesize 10--14 & \footnotesize 2 & \footnotesize 10 & \CheckBox[name=ok_amoxicillin_57,width=1.6ex,height=1.6ex]{}\\
\footnotesize 14--19 & \footnotesize 3 & \footnotesize 15 & \CheckBox[name=ok_amoxicillin_58,width=1.6ex,height=1.6ex]{}\\
\end{longtable}
\subsection*{cotrimoxazole\_prophylaxis}
{\footnotesize prophylaxis in HIV confirmed or exposed child \textperiodcentered{} oral \textperiodcentered{} once a day, starting at 4-6 weeks of age \textperiodcentered{} source 2014 p16 \textperiodcentered{} keyed by age}\\[-2pt]
\begin{longtable}{@{}l r r r c@{}}
\rowcolor{rule!40}\footnotesize\textbf{age (months)} & \footnotesize\textbf{syrup\_40\_200\_per\_5ml\_ml} & \footnotesize\textbf{paed\_tablet\_20\_100} & \footnotesize\textbf{adult\_tablet\_80\_400} & \footnotesize\textbf{OK}\\
\endhead
\footnotesize 0--6 & \footnotesize 2.5 & \footnotesize 0.5 & \footnotesize 0.25 & \CheckBox[name=ok_cotrimoxazole_prophylaxis_59,width=1.6ex,height=1.6ex]{}\\
\footnotesize 6--60 & \footnotesize 5 & \footnotesize 2 & \footnotesize 0.5 & \CheckBox[name=ok_cotrimoxazole_prophylaxis_60,width=1.6ex,height=1.6ex]{}\\
\end{longtable}
\subsection*{paracetamol}
{\footnotesize high fever (>38.5C) or ear pain \textperiodcentered{} oral \textperiodcentered{} every 6 hours until fever or pain is gone \textperiodcentered{} source 2014 p16 \textperiodcentered{} keyed by weight}\\[-2pt]
\begin{longtable}{@{}l r r c@{}}
\rowcolor{rule!40}\footnotesize\textbf{weight (kg)} & \footnotesize\textbf{tablet\_100mg} & \footnotesize\textbf{tablet\_500mg} & \footnotesize\textbf{OK}\\
\endhead
\footnotesize 4--14 & \footnotesize 1 & \footnotesize 0.25 & \CheckBox[name=ok_paracetamol_61,width=1.6ex,height=1.6ex]{}\\
\footnotesize 14--19 & \footnotesize 1.5 & \footnotesize 0.5 & \CheckBox[name=ok_paracetamol_62,width=1.6ex,height=1.6ex]{}\\
\end{longtable}
\subsection*{zinc}
{\footnotesize diarrhoea (age 2 months up to 5 years) \textperiodcentered{} oral \textperiodcentered{} daily for 14 days \textperiodcentered{} source 2014 (GIVE EXTRA FLUID FOR DIARRHOEA) \textperiodcentered{} keyed by age}\\[-2pt]
\begin{longtable}{@{}l r c@{}}
\rowcolor{rule!40}\footnotesize\textbf{age (months)} & \footnotesize\textbf{tablets\_daily} & \footnotesize\textbf{OK}\\
\endhead
\footnotesize 2--6 & \footnotesize 0.5 & \CheckBox[name=ok_zinc_63,width=1.6ex,height=1.6ex]{}\\
\footnotesize 6--$\infty$ & \footnotesize 1 & \CheckBox[name=ok_zinc_64,width=1.6ex,height=1.6ex]{}\\
\end{longtable}
\subsection*{ors\_plan\_a\_extra\_fluid}
{\footnotesize diarrhoea, no dehydration -- amount after each loose stool \textperiodcentered{} oral \textperiodcentered{}  \textperiodcentered{} source 2014 Plan A \textperiodcentered{} keyed by age}\\[-2pt]
\begin{longtable}{@{}l r r c@{}}
\rowcolor{rule!40}\footnotesize\textbf{age (months)} & \footnotesize\textbf{ml\_after\_each\_loose\_stool\_min} & \footnotesize\textbf{ml\_after\_each\_loose\_stool\_max} & \footnotesize\textbf{OK}\\
\endhead
\footnotesize 0--24 & \footnotesize 50 & \footnotesize 100 & \CheckBox[name=ok_ors_plan_a_extra_fluid_65,width=1.6ex,height=1.6ex]{}\\
\footnotesize 24--$\infty$ & \footnotesize 100 & \footnotesize 200 & \CheckBox[name=ok_ors_plan_a_extra_fluid_66,width=1.6ex,height=1.6ex]{}\\
\end{longtable}
\subsection*{vitamin\_a}
{\footnotesize supplementation from 6 months; treatment dose for measles or persistent diarrhoea (NOT if dosed in past month or on RUTF) \textperiodcentered{} oral \textperiodcentered{}  \textperiodcentered{} source 2014 (Vitamin A supplementation and treatment) \textperiodcentered{} keyed by age}\\[-2pt]
\begin{longtable}{@{}l r c@{}}
\rowcolor{rule!40}\footnotesize\textbf{age (months)} & \footnotesize\textbf{iu} & \footnotesize\textbf{OK}\\
\endhead
\footnotesize 6--12 & \footnotesize 100000 & \CheckBox[name=ok_vitamin_a_67,width=1.6ex,height=1.6ex]{}\\
\footnotesize 12--$\infty$ & \footnotesize 200000 & \CheckBox[name=ok_vitamin_a_68,width=1.6ex,height=1.6ex]{}\\
\end{longtable}
\subsection*{mebendazole}
{\footnotesize deworming if hookworm/whipworm endemic, child >=1 year, no dose in previous 6 months \textperiodcentered{} oral \textperiodcentered{} single dose in clinic \textperiodcentered{} source 2014 \textperiodcentered{} keyed by age}\\[-2pt]
\begin{longtable}{@{}l r c@{}}
\rowcolor{rule!40}\footnotesize\textbf{age (months)} & \footnotesize\textbf{dose\_mg} & \footnotesize\textbf{OK}\\
\endhead
\footnotesize 12--$\infty$ & \footnotesize 500 & \CheckBox[name=ok_mebendazole_69,width=1.6ex,height=1.6ex]{}\\
\end{longtable}
\subsection*{ceftriaxone\_im}
{\footnotesize sick-child pre-referral (very severe disease, suspected meningitis, severe pneumonia, mastoiditis) \textperiodcentered{} IM \textperiodcentered{}  \textperiodcentered{} source 2022 p35 \textperiodcentered{} keyed by weight}\\[-2pt]
{\footnotesize note: >17.5 kg: give 2 ml in each thigh; for weights over 17.5 kg dilute 1 g in 3.5 ml sterile water and give 5.5 ml IM}\\[-2pt]
\begin{longtable}{@{}l r r c@{}}
\rowcolor{rule!40}\footnotesize\textbf{weight (kg)} & \footnotesize\textbf{dose\_mg} & \footnotesize\textbf{volume\_ml} & \footnotesize\textbf{OK}\\
\endhead
\footnotesize 3.5--5.5 & \footnotesize 312 & \footnotesize 1.25 & \CheckBox[name=ok_ceftriaxone_im_70,width=1.6ex,height=1.6ex]{}\\
\footnotesize 5.5--7 & \footnotesize 440 & \footnotesize 1.75 & \CheckBox[name=ok_ceftriaxone_im_71,width=1.6ex,height=1.6ex]{}\\
\footnotesize 7--9 & \footnotesize 625 & \footnotesize 2.5 & \CheckBox[name=ok_ceftriaxone_im_72,width=1.6ex,height=1.6ex]{}\\
\footnotesize 9--11 & \footnotesize 750 & \footnotesize 3 & \CheckBox[name=ok_ceftriaxone_im_73,width=1.6ex,height=1.6ex]{}\\
\footnotesize 11--14 & \footnotesize 810 & \footnotesize 3.25 & \CheckBox[name=ok_ceftriaxone_im_74,width=1.6ex,height=1.6ex]{}\\
\footnotesize 14--17.5 & \footnotesize 1000 & \footnotesize 4 & \CheckBox[name=ok_ceftriaxone_im_75,width=1.6ex,height=1.6ex]{}\\
\footnotesize 17.5--$\infty$ & \footnotesize 1500 & \footnotesize 5.5 & \CheckBox[name=ok_ceftriaxone_im_76,width=1.6ex,height=1.6ex]{}\\
\end{longtable}
\subsection*{diazepam\_rectal}
{\footnotesize convulsing now \textperiodcentered{} rectal \textperiodcentered{}  \textperiodcentered{} source 2022 p35 \textperiodcentered{} keyed by weight}\\[-2pt]
{\footnotesize note: 0.5 mg/kg per rectum; repeat after 10 minutes if not stopped}\\[-2pt]
\begin{longtable}{@{}l r r c@{}}
\rowcolor{rule!40}\footnotesize\textbf{weight (kg)} & \footnotesize\textbf{dose\_mg} & \footnotesize\textbf{volume\_ml} & \footnotesize\textbf{OK}\\
\endhead
\footnotesize 3--4 & \footnotesize 2 & \footnotesize 0.4 & \CheckBox[name=ok_diazepam_rectal_77,width=1.6ex,height=1.6ex]{}\\
\footnotesize 4--5 & \footnotesize 2.5 & \footnotesize 0.5 & \CheckBox[name=ok_diazepam_rectal_78,width=1.6ex,height=1.6ex]{}\\
\footnotesize 5--15 & \footnotesize 5 & \footnotesize 1 & \CheckBox[name=ok_diazepam_rectal_79,width=1.6ex,height=1.6ex]{}\\
\footnotesize 15--25 & \footnotesize 7.5 & \footnotesize 1.5 & \CheckBox[name=ok_diazepam_rectal_80,width=1.6ex,height=1.6ex]{}\\
\end{longtable}
\subsection*{young\_infant\_ampicillin\_im}
{\footnotesize young infant (0-2 months) very severe disease -- with gentamicin \textperiodcentered{} IM \textperiodcentered{}  \textperiodcentered{} source 2014 p51 \textperiodcentered{} keyed by weight}\\[-2pt]
{\footnotesize note: 2014 young-infant regimen. Dose 50 mg/kg. 250 mg vial + 1.3 ml sterile water = 250 mg/1.5 ml.}\\[-2pt]
\begin{longtable}{@{}l r c@{}}
\rowcolor{rule!40}\footnotesize\textbf{weight (kg)} & \footnotesize\textbf{volume\_ml} & \footnotesize\textbf{OK}\\
\endhead
\footnotesize 1--1.5 & \footnotesize 0.4 & \CheckBox[name=ok_young_infant_ampicillin_im_81,width=1.6ex,height=1.6ex]{}\\
\footnotesize 1.5--2 & \footnotesize 0.5 & \CheckBox[name=ok_young_infant_ampicillin_im_82,width=1.6ex,height=1.6ex]{}\\
\footnotesize 2--2.5 & \footnotesize 0.7 & \CheckBox[name=ok_young_infant_ampicillin_im_83,width=1.6ex,height=1.6ex]{}\\
\footnotesize 2.5--3 & \footnotesize 0.8 & \CheckBox[name=ok_young_infant_ampicillin_im_84,width=1.6ex,height=1.6ex]{}\\
\footnotesize 3--3.5 & \footnotesize 1.0 & \CheckBox[name=ok_young_infant_ampicillin_im_85,width=1.6ex,height=1.6ex]{}\\
\footnotesize 3.5--4 & \footnotesize 1.1 & \CheckBox[name=ok_young_infant_ampicillin_im_86,width=1.6ex,height=1.6ex]{}\\
\footnotesize 4--4.5 & \footnotesize 1.3 & \CheckBox[name=ok_young_infant_ampicillin_im_87,width=1.6ex,height=1.6ex]{}\\
\end{longtable}
\subsection*{artemether\_lumefantrine}
{\footnotesize uncomplicated malaria (malaria test positive) \textperiodcentered{} oral \textperiodcentered{} two times daily for 3 days \textperiodcentered{} source 2014 p16 \textperiodcentered{} keyed by weight}\\[-2pt]
{\footnotesize note: first dose in clinic, observe 1 hour, repeat if vomits within an hour; take with food}\\[-2pt]
\begin{longtable}{@{}l r c@{}}
\rowcolor{rule!40}\footnotesize\textbf{weight (kg)} & \footnotesize\textbf{tablets\_per\_dose} & \footnotesize\textbf{OK}\\
\endhead
\footnotesize 5--10 & \footnotesize 1 & \CheckBox[name=ok_artemether_lumefantrine_88,width=1.6ex,height=1.6ex]{}\\
\footnotesize 10--14 & \footnotesize 1 & \CheckBox[name=ok_artemether_lumefantrine_89,width=1.6ex,height=1.6ex]{}\\
\footnotesize 14--19 & \footnotesize 2 & \CheckBox[name=ok_artemether_lumefantrine_90,width=1.6ex,height=1.6ex]{}\\
\end{longtable}
\section*{Specific things to check}
\begin{itemize}[leftmargin=2.4em,itemsep=5pt,label={}]
\item \CheckBox[name=ok_chk_91,width=1.6ex,height=1.6ex]{}\hspace{0.4em}\footnotesize Confirm every weight/age band boundary against the source table (off-by-one on a band edge is a dosing error).
\item \CheckBox[name=ok_chk_92,width=1.6ex,height=1.6ex]{}\hspace{0.4em}\footnotesize Ceftriaxone dilution and per-thigh split (>17.5kg) -- 2022 p35.
\item \CheckBox[name=ok_chk_93,width=1.6ex,height=1.6ex]{}\hspace{0.4em}\footnotesize 2014 uses ampicillin+gentamicin for the young infant; 2022 sick-child pre-referral is ceftriaxone. Confirm which the national adaptation wants.
\item \CheckBox[name=ok_chk_94,width=1.6ex,height=1.6ex]{}\hspace{0.4em}\footnotesize ART regimens are edition- and country-specific -- verify against current national ART guidelines, not this booklet.
\item \CheckBox[name=ok_chk_95,width=1.6ex,height=1.6ex]{}\hspace{0.4em}\footnotesize Zinc, ORS volumes, vitamin A, mebendazole are stable across editions but confirm.
\item \CheckBox[name=ok_chk_96,width=1.6ex,height=1.6ex]{}\hspace{0.4em}\footnotesize Fever severe label kept as very\_severe\_febrile\_disease (2014) with bulging fontanelle added -- acceptable?
\item \CheckBox[name=ok_chk_97,width=1.6ex,height=1.6ex]{}\hspace{0.4em}\footnotesize HIV severity: all tiers MODERATE except hiv\_infection\_unlikely (MILD), none PINK (ART not urgent) -- acceptable?
\item \CheckBox[name=ok_chk_98,width=1.6ex,height=1.6ex]{}\hspace{0.4em}\footnotesize Young-infant treatment emits no specific dose (IM antibiotics / referral only) -- acceptable?
\end{itemize}
{\footnotesize Full per-branch checklist items are in the \texttt{src/sft/extended\_protocol.py} and \texttt{src/sft/young\_infant.py} docstrings.}\\[10pt]
\section*{Sign-off}
\begin{tabular}{@{}p{5.5cm}p{5.5cm}p{4.5cm}@{}}
\TextField[name=reviewer,width=5.2cm]{} & \TextField[name=signature,width=5.2cm]{} & \TextField[name=date,width=4.2cm]{}\\
\footnotesize Reviewer name & \footnotesize Signature & \footnotesize Date\\
\end{tabular}\\[8pt]
\TextField[name=adaptation,width=\linewidth]{}\\[-2pt]
{\footnotesize National adaptation reviewed against}
\end{Form}
\end{document}